\documentclass[11pt,a4paper]{article}
\usepackage[UTF8]{ctex}
\usepackage[margin=2.2cm]{geometry}
\usepackage{booktabs}
\usepackage{array}
\usepackage{tabularx}
\usepackage{enumitem}
\usepackage{hyperref}
\hypersetup{colorlinks=true,linkcolor=blue,urlcolor=blue}
\newcommand{\RepoURL}{https://github.com/itachiliu/Energy-Aware-Deep-Learning-Compilation}
\title{Supplementary Material\\[4pt]\large 能效感知深度学习编译综述：编码、证据与可复现性}
\author{何欣程、刘岩}
\date{2026-09-23}
\begin{document}
\maketitle

\begin{center}
\fbox{\parbox{0.94\textwidth}{\small
\textbf{本文件的机器可读版本与全部数据文件见配套仓库：}\\
\url{\RepoURL}\\
正文中出现的“补充材料”均指该仓库。}}
\end{center}

\section*{A. 本文件与文件索引}
正文之外的所有证据按用途分成七组。文件路径相对仓库根目录：

\begin{itemize}[leftmargin=2em,itemsep=1pt]
  \item \textbf{S1 · 0/18 证据}：\texttt{ecosystem/ecosystem-selection.md}；
  \item \textbf{S2 · 矩阵计数}：\texttt{matrix/matrix-counts-4x4.md}、
    \texttt{matrix-counts-4x4.tex}、\texttt{matrix/count\_matrix.py}；
  \item \textbf{S3 · 编码与一致性}：\texttt{coding/coding-121-final.csv}、
    \texttt{coding-sheet-A/B-v2.csv}、\texttt{-v3.csv}、
    \texttt{codebook-*-worked-examples.csv}、
    \texttt{reconciliation-sheet.csv}、
    \texttt{reconciliation-resolution-1-full.csv}、
    \texttt{reconciliation-resolution-2.csv}、
    \texttt{reconciliation-resolution-final.txt}、
    以及 \texttt{coding/kappa/} 下的 $\kappa$ 脚本与报告；
  \item \textbf{S4 · 历史审计}：\texttt{audit/appendix-matrix-code.csv}、
    \texttt{code\_matrix.py}、\texttt{audit\_internal.py}、
    \texttt{internal-audit.md}；
  \item \textbf{S5 · 独立抽查}：\texttt{audit/coding-audit-sample30.csv}、
    \texttt{coding-audit-protocol.md}、\texttt{audit\_coding.py}；
  \item \textbf{S6 · 最小验证实验}：\texttt{experiment/}（协议、原始数据、
    作业包与绘图脚本）；
  \item \textbf{S7 · 追踪与环境}：\texttt{methodology/2026-preprint-tracking.md}、
    \texttt{experiment/ENV\_REQUIREMENTS.md}；
  \item \textbf{检索与筛选}：\texttt{methodology/prisma\_counts.csv}，
    逐库重建的复核用检索式见论文附录~B。
\end{itemize}

\section{S1 · 0/18 生态表：系统与证据载体}
打标依据优先级：论文 $>$ 官方文档 $>$ 公开源码；标注源码位置的条目可直接复核。
0/18 的范围为“本综述选取的 18 个代表性通用/生产级栈（16 个独立项目）”，
时间切片为 2026-05。

{\scriptsize
\begin{tabularx}{\textwidth}{@{}clXX@{}}
\toprule
\# & 系统/组件 & 论文/文档证据 & 源码/可复核点 \\
\midrule
1 & TVM & chen\_tvm\_2018 & 开源；ProgramMeasurer 位于 \texttt{src/auto\_scheduler/measure.cc} \\
2 & Ansor（TVM） & zheng\_ansor\_2020 & 同上；代价模型以延迟为默认预测目标 \\
3 & TVM Unity / Relax & tvm\_unity & 官方仓库/文档（版本定稿固定） \\
4 & TorchInductor & ansel\_pytorch\_2024 & 官方仓库/文档（版本定稿固定） \\
5 & Triton & triton & 官方仓库（版本定稿固定） \\
6 & XLA & openxla & ProfileGuidedLatencyEstimator；GPU Speed-of-Light 代价模型 \\
7 & IREE & iree & 官方仓库/文档（版本定稿固定） \\
8 & TensorRT & nvidia\_tensorrt & 文档级；闭源，外围组件开源 \\
9 & TensorRT-LLM & nvidia\_trtllm & 开源仓库；执行依赖闭源 TRT 运行时 \\
10 & OpenVINO & intel\_openvino & 官方仓库/文档（版本定稿固定） \\
11 & ONNX Runtime & onnxruntime & 官方仓库/文档（版本定稿固定） \\
12 & TFLite / NNAPI & google\_tflite & 文档级；ANEURALNETWORKS\_PREFER\_LOW\_POWER \\
13 & Core ML & apple\_coreml & 文档级；MLComputeUnits 枚举 \\
14 & ExecuTorch & executorch & 官方仓库/文档（版本定稿固定） \\
15 & MNN & mnn & 官方仓库（版本定稿固定） \\
16 & NCNN & ncnn & 官方仓库（版本定稿固定） \\
17 & Paddle Lite & paddlelite & 官方仓库（版本定稿固定） \\
18 & MLC-LLM & mlcllm & 官方仓库（版本定稿固定） \\
\bottomrule
\end{tabularx}}

候选未入选的正文可追溯项：vLLM 等服务期连续批处理运行时（决策主体为服务期）、
2026 预印本（MileStone/AutoPass/DualScale）、torch.compile 的 Dynamo 前端
（后端由 TorchInductor 承担）。

\section{S2 · 4×4 矩阵计数与空格}
计数口径：仅统计粒度与机制均落在四类中的附录 A 样本（$n=68$）；编码取自
\texttt{coding-121-final.csv}，空格判定结合附录编码与正文代表
（计数脚本 \texttt{matrix/count\_matrix.py}）。

\begin{center}
\begin{tabular}{lcccc}
\toprule
粒度单元 & Analytical & Learned & AutoSearch & Agentic \\
\midrule
计算图（Graph） & 6 & 3 & 11 & 0 \\
循环/算子（Loop/Op） & 4 & 4 & 19 & 1 \\
异构放置（Placement） & 2 & 3 & 5 & 0 \\
功耗状态（Power） & 4 & 3 & 3 & 0 \\
\bottomrule
\end{tabular}
\end{center}

行合计：Graph 20、Loop/Op 28、Placement 10、Power 10。

空格（附录计数为 0 且正文无代表）：Placement × Agentic、Power × Agentic。
Graph × Agentic 计数为 0，但以正文非统计代表（Agentic MLIR）视为非空格。

\section{S3 · 两位编码者独立编码与一致性}
编码字段为类别、粒度、机制、对象类型四项，均由两位编码者按公开受控词表独立完成；
编码手册与示例见 \texttt{coding/coding-rubric-v2.md}、\texttt{coding-rubric-v3.md} 与
\texttt{codebook-v2-worked-examples.csv}、\texttt{codebook-v3-worked-examples.csv}。
示例表所固定的单元不计入一致性统计。

\begin{center}
\begin{tabular}{lrrr}
\toprule
字段 & $N$ & 一致率 & Cohen's $\kappa$ \\
\midrule
类别 & 118 & 79.7\% & 0.77 \\
粒度 & 118 & 74.6\% & 0.66 \\
机制 & 114 & 73.7\% & 0.65 \\
对象类型 & 118 & 89.8\% & 0.85 \\
\bottomrule
\end{tabular}
\end{center}

机制与对象类型字段在编码手册扩充后进行了第二轮独立编码；第一轮机制的
$\kappa=0.56$，其偏低的取值正是扩充手册中 M1/M2 判据并重做该轮的动因。
早期版本中的“核心/桥接/前沿”边界标签 $\kappa=0.46$，已由对象类型字段取代。
完整脚本与报告见 \texttt{coding/kappa/}。

独立编码共产生 96 个不一致单元（类别 24、粒度 30、机制 30、对象类型 12）。
两位编码者先各自逐条复议一遍，31 个单元收敛；对其余 65 个单元改按“先定判据、
再落单元”的方式完成第三轮联合复核，形成覆盖全部 96 个单元的决议表
（\texttt{reconciliation-resolution-final.txt}）：48 个采编码者 A 的取值、
16 个采编码者 B 的取值、1 个采用第三值（S01 的粒度）；一度标注待核或注记
冲突的 3 个单元（S18、S19、S68 的机制）已回查原文后确认。争议清单见
\texttt{reconciliation-sheet.csv}，前两轮复议记录见
\texttt{reconciliation-resolution-1-full.csv} 与
\texttt{reconciliation-resolution-2.csv}，第三轮裁定包见
\texttt{adjudication-round3-guide.md}，归并脚本见
\texttt{coding/kappa/merge\_reconciliation.py}。

\section{S4 · 121 行编码文件与历史审计}
\begin{itemize}[leftmargin=2em,itemsep=1pt]
  \item 121 行全部给出类别、粒度、机制、对象类型四列编码
        （\texttt{coding-121-final.csv}；附录镜像
        \texttt{audit/appendix-matrix-code.csv}）；
  \item 单编码者阶段的规则式编码与差异清单（\texttt{audit/code\_matrix.py}、
        \texttt{internal-audit.md}）保留作历史留痕：在该阶段，121 行中 23 行由
        S 级覆盖改变默认规则；两人独立编码取代它之后，同一规则与最终编码的
        差异行数升至 87，说明规则式编码无法替代人工判定；
  \item 六处易争议条目（WELDER、HAQ、Checkmate、Stream、ODiMO、PowerFlow-DNN）
        的正文依据见 \texttt{matrix/matrix-empty-cells.md}。
\end{itemize}

\section{S5 · 独立抽查协议（可选）}
若期刊或审稿人要求非作者抽查：由第三方按
\texttt{audit/coding-audit-protocol.md} 填写 \texttt{coding-audit-sample30.csv}
的独立列，运行 \texttt{audit/audit\_coding.py} 输出一致率与 Cohen's $\kappa$。

\section{S6 · 最小验证实验}
本文“小规模验证实验”一节报告的实验，其测量协议、原始逐次测量、配置矩阵、
作业包与绘图脚本均在 \texttt{experiment/} 下：说明见
\texttt{experiment/README.md}，环境与验收命令见
\texttt{experiment/ENV\_AND\_EXPERIMENT.md}，逐次测量见
\texttt{experiment/data/}，作业包见 \texttt{experiment/job/}，
绘图脚本见 \texttt{experiment/scripts/}。

\section{S7 · 2026 窗口后追踪与复现环境}
2026 年 5 月检索窗口后的预印本追踪（MileStone、AutoPass、DualScale）清单见
\texttt{methodology/2026-preprint-tracking.md}；批次二实验所需环境与验收命令见
\texttt{experiment/ENV\_REQUIREMENTS.md}。

\section*{B. 版本与验证状态}
\begin{itemize}[leftmargin=2em,itemsep=1pt]
  \item 正文：XeLaTeX 连编通过，无未定义引用；
  \item 编码：\texttt{coding/kappa/merge\_reconciliation.py}、
    \texttt{matrix/count\_matrix.py} 与
    \texttt{coding/kappa/sync\_appendix\_from\_final.py} 可复跑；
  \item 一致性审计：\texttt{audit/audit\_scope\_appendix\_consistency.py} 报告
    \texttt{no inconsistencies found}；\texttt{audit/crosscheck\_appendix.py}
    报告样本表与机器副本 0 处不匹配；
  \item 编码者：两位编码者独立编码，经两轮复议与第三轮联合裁定；
    尚未引入非作者独立抽查（S5 为可选补强项）。
\end{itemize}

\end{document}
